\documentclass[10pt,a4paper]{amsart}
\usepackage[margin=19mm]{geometry}
\usepackage{amsmath,amssymb,mathtools}
\usepackage[scr=rsfs,cal=euler]{mathalfa}
\usepackage{xfrac}
\usepackage[unicode,hidelinks,bookmarks=false]{hyperref}
\newcommand{\ie}{i.e.\ }
\newcommand{\cf}{cf.\ }
\newcommand{\resp}{respectively\ }
\newcommand{\Subsection}{\subsection}
\providecommand{\nobreakdash}{\nobreak\hbox{-}}
\setlength{\emergencystretch}{2em}
\multlinegap0pt
\usepackage{amssymb}
\usepackage{mathtools}
\usepackage{bm}
\usepackage[shortlabels]{enumitem}
\usepackage{tikz}
\newtheorem{theorem}{Theorem}
\newtheorem{corollary}{Corollary}
\usepackage{cleveref}
\crefname{theorem}{Theorem}{Theorems}
\crefname{corollary}{Corollary}{Corollarys}
\newcommand{\Z}{\mathbb Z}
\newcommand{\cF}{\mathcal F}
\title{A Conway--Coxeter theorem \\for decorated frieze patterns}
\author{Takeru Kuwana, Katsuhiko Matsuzaki}
\date{Source: arXiv:2609.06029v1 (CC BY 4.0)}
\begin{document}
\maketitle

\noindent\emph{Source and licence.} A Conway--Coxeter theorem \\for decorated frieze patterns. Version arXiv:2609.06029v1 (CC BY 4.0), distributed by arXiv under the Creative Commons Attribution 4.0 International licence, http://creativecommons.org/licenses/by/4.0/, as recorded in the arXiv licence metadata for this exact version and shown on the abstract page https://arxiv.org/abs/2609.06029v1. Full licence notice: \texttt{licenses/arxiv-2609.06029-CC-BY-4.0.txt}.\par\smallskip

\noindent\textit{Editorial scope.} Counted unit: the article's corollary eq:fundamental-domain (source0.tex lines 828--834). The article's complete original proof of it is delivered with it (source0.tex lines 836--842, one \texttt{proof} environment of 7 lines). No mathematics is written, repaired or re-derived: the statement and the proof are the article's own lines, in the article's own notation, and no retained line is substituted or re-flowed.  Retained uncounted as the source-local context the proof needs: lines 563--569; lines 698--700.  Nothing was omitted from any retained span, so the package carries no omission record.  No retained line contains a citation, so the wrapper carries no bibliography and no bibliography entry.  No retained line contains a figure environment, an image-inclusion command or drawing code, so no image is delivered and the package declares no asset dependency.  Editorial formatting only, in addition: an AMS \texttt{amsart} preamble that restates the article's own declarations and macros used by the retained lines, namely the environments \texttt{theorem}, \texttt{corollary} on separate counters, together with the standard packages those lines need.  The retained environments print in this document as Theorem 1, Corollary 1 in order of appearance, whereas the article prints the same items with the numbering of its own full document.  The complete native source of the article as distributed in the arXiv e-print for this exact version is bundled unchanged as \texttt{sources/209485/source0.tex}.
\begin{theorem}[Algebraic periodicity]\label{thm:periodicity}
Let $\cF$ be a closed decorated frieze of width $m$, and put $n=m+3$. Then
$a_{i+n}=a_i$ $(i\in\Z)$,
and consequently
$W_{i+n,j+n}=W_{i,j}$
throughout the strip. Thus $n$ is a period of every closed decorated frieze of width $m$.
\end{theorem}

\begin{equation}\label{eq:glide}
 W_{i,j}=W_{j,i+n}
\end{equation}

\begin{corollary}[Triangular fundamental domain]\label[corollary]{cor:fundamental-domain}
Every closed decorated frieze of width $m$ is completely determined by the finite triangular collection
\begin{equation*}\label{eq:fundamental-domain}
 \mathscr D_n=\{W_{i,j}:1\le i<j\le n\}.
\end{equation*}
Indeed, every nonzero entry of the infinite strip is carried to a unique entry of $\mathscr D_n$ by iterating the glide reflection $G$ (equivalently, by using $n$-periodicity together with \eqref{eq:glide}).
\end{corollary}

\begin{proof}
Let $W_{i,j}$ be nonzero, so $1\le j-i\le n-1$. By \cref{thm:periodicity}, translate both indices by a multiple of $n$ so that $1\le i\le n$. If then $j\le n$, the entry belongs to $\mathscr D_n$. If $j>n$, the glide relation gives
\[
 W_{i,j}=W_{j,i+n}=W_{j-n,i},
\]
where the last equality uses $n$-periodicity. Since $1\le j-n<i\le n$, this lies in $\mathscr D_n$. The reverse reconstruction is immediate from periodicity and the glide relation. Uniqueness of the representative follows because, modulo the translation $G^2(i,j)=(i+n,j+n)$, a $G$-orbit is determined by the unordered pair of residue classes $\{i,j\}$ modulo $n$, and $\mathscr D_n$ contains exactly one ordered representative $1\le i<j\le n$ of each such pair.
\end{proof}

\end{document}
